\scanpage{48}
\mspage{4)}
On a un foncteur
\diagram{\begin{tikzcd}[ampersand replacement=\&]
(*)\quad\Gamma:\quad S_*\mathcal TD^\pm_{\mathrm{top}}\arrow[r]\&
\deleted{\unsure{\mathrm{F}}\mathrm{MD}}(\mathrm{Gra\,bipond})
\end{tikzcd}}
qui associe à un objet \inserted{$(X,\widetilde\Sigma,R)$} de
$S_*\mathcal TD^\pm_{\mathrm{top}}$ le graphe \inserted{MD}
$(S,\vec A,\sigma,\deleted{\unsure{g}}\inserted{\unsure{o|\sigma}},\deleted{\unsure{)}}\gamma,\rho)$ suivant :
\[
S=\pi_0(X\deleted{\unsure{)|,}}\setminus\widetilde\Sigma)\xleftarrow{\sim}\pi_0(\text{Déc}(X,\widetilde\Sigma))
\]
$\vec A=$ ens des couples $(C,\deleted{\unsure{\xi}}\omega)$, où
$C$ \deleted{\unsure{$\in A$|est un}}\inserted{$\in\pi_0(\widetilde\Sigma)$}, et où
$\omega\in\mathrm{Or}(C)\inserted{\simeq\mathrm{rive}_X(C)}$,
soumis à la condition que si $C\in\pi_0(\partial X)$, \unsure{alors|donc} $\omega$
soit l'orientation induite i.e. la rive qui se trouve dans $X$

$\sigma=$ échange des rives (quand on a le droit d'échanger,
sinon, i.e. si $C\in\pi_0(\partial X)$, on pose $\sigma(a)=a$)

$\deleted{\unsure{O}}\inserted{o}$ : associe à $(C,\text{rive de }C)$
l'unique ``face'' $\in\pi_0(X\setminus\widetilde\Sigma)$ qui
\deleted{$\unsure{R\subset\widetilde R}$} contient cette rive
\diagram{\begin{tikzcd}[ampersand replacement=\&]
\gamma:\quad S\arrow[r]\arrow[d,equal]\&\mathbb N\\
\pi_0(X\setminus\widetilde\Sigma)\&
\end{tikzcd}}
application \uline{genre} des composants $X_i$.
\diagram{\begin{tikzcd}[ampersand replacement=\&]
\deleted{\unsure{\rho}}\inserted{r}:\quad A=\pi_0(\widetilde\Sigma)\arrow[r]\&\mathbb N\&
:\quad C\arrow[r,mapsto]\&\mathrm{Card}(R\cap C).
\end{tikzcd}}

Le foncteur $(*)$ \deleted{est \uline{ess. surjectif} et}
\inserted{induit un \uline{iso} sur les} $\pi_0$ et une surjectif
sur les $\pi_1$ (c'est un \uline{0-isomorphisme} \deleted{\unsure{)|?}},
en langage homotopique).

\deleted{\unsure{Prop.|La preuve}} :

\uline{Remarques}. Dans $S_*\mathcal TD^\pm_{\mathrm{top}}$ il y a une
opération évidente de \uline{somme} \inserted{\uline{directe}},
dans $\deleted{\mathrm{FMD}}\inserted{\mathrm{Gra\,bipond}}$ aussi,
et $(*)$ y commute.
\mspage{47}
